\documentclass
[
	12pt,
	english,
	listof=totoc,
	bibliography=totoc,
	parskip,
%	numbers=endperiod
]{scrartcl}

\usepackage{geometry}
\usepackage[english]{babel}
\usepackage{amsmath}
\usepackage{amsfonts}
\usepackage{amssymb}
\usepackage{booktabs}					% Nicer tables
\setlength{\abovetopsep}{1ex}

\usepackage{graphicx}
\usepackage{xcolor}
\usepackage[utf8]{inputenc}		% Input encoding
\usepackage{csquotes}					% Quotations
\usepackage{varioref}					% References like "Figure 2 on the next page"
\usepackage[colorlinks=true, allcolors=blue!50!black]{hyperref}
\usepackage[T1]{fontenc}			% Output encoding
\usepackage{libertinus}				% Load Libertinus fonts
\usepackage{lmodern}
\usepackage{booktabs}

%\setmonofont{DejaVu Sans Mono}[Scale=MatchLowercase]

%\setuptoc{toc}{totoc} % Add 'Contents' to the table of contents
\setcounter{tocdepth}{2} % suppress subsubsections

\usepackage[headsepline]{scrlayer-scrpage}	% Horizontal line under header
\clearpairofpagestyles
\automark{section}
\ohead{\normalfont\headmark}
\ofoot{\normalfont\thepage}

\clubpenalty=10000
\widowpenalty=10000
\displaywidowpenalty=10000

\usepackage[capitalise,noabbrev,nameinlink]{cleveref}	% Automatically add type names like "Figure"

\usepackage[
		backend=biber,
		style=numeric,
		sorting=nyt,
		backref=true
]{biblatex}
\addbibresource{references.bib}

\author{Michael Heep}

\begin{document}

\begin{titlepage}
    \begin{minipage}[c]{0.45\textwidth}
    \includegraphics[width=0.625\textwidth]{resources/UNIS-Dachmarke_cmyk.pdf}
    \end{minipage}\hfill
    \begin{minipage}[c]{0.45\textwidth}
    \includegraphics[width=1\textwidth]{resources/UNIS-f4-nt_cmyk.pdf}
    \end{minipage}	
	\vspace*{1cm}
    \begin{center}
        \vspace*{2cm}
        \huge
	Master's Thesis \\
        \vspace*{1cm}
	% Version 1:Agentic AI for CrypTool-Online Developers\\
  % Version 2 DE (Wismüller): Konzepte zur produktiven Nutzung von LLMs (AI Agents) in der Entwicklung von großen Softwareprojekten am Beispiel von CrypTool\\
  % Version 2 EN (Wismüller): Concepts for the Productive Use of LLMs (AI Agents) in the Development of Large Software Projects with CrypTool as an Example
  % Version 3 DE (BE+MH): Evaluation und Verbesserung des produktiven Entwicklungsprozesses von großen Softwareprojekten mit Hilfe von LLMs (AI Agents) am Beispiel von CrypTool\\
  % Version 3 EN (BE+MH): Evaluation and Improvement of the Productive Development Process of Large Software Projects with the Help of LLMs (AI Agents) with CrypTool as an Example\\
  % Version 4 EN:Evaluation and Improvement of LLMs (AI Agents) Contribution to Large Software Projects with CrypTool as an Example\\
% Final Version:
  Evaluating and Improving Software Development Processes in Large-Scale Projects Using LLM-Based AI Agents: A Case Study of CrypTool
        
  
    \end{center}
    \vfill
    \normalsize
    \begin{tabular}{ll}
	Examiner & Prof. Dr. Bernhard Esslinger \\
	& Prof. Dr. Roland Wismüller \\
        %Supervisor & % TODO \\
        \\
        Author & Michael Heep \\
        Submitted & % TODO
    \end{tabular}
\end{titlepage}

%%%
%%% Kurzzusammenfassung
%%%
\selectlanguage{german}
\addsec{Kurzzusammenfassung}

% TODO: Write summary in German

%\newpage
\selectlanguage{english}
%%%
%%% Abstract (kann auch auf dieselbe Seite wie die Zusammenfassung)
%%%
\addsec{Abstract}
% TODO: Write summary in English


%%% 
%%% Table of Contents
%%%
\newpage
\tableofcontents
\newpage

%%%
%%% Introduction section
%%% Motivation, problem statement, research questions, delimitations and structure of the thesis
%%%
\section{Introduction} % TODO: Put methodology somewhere into introduction

\subsection{Motivation and Problem Statement} 
% Companies try to use LLMs (AI Agents) to accelerate their software development process, but there is a lack of concepts and best practices for how to do this effectively
% Developers who are not familiar with a new software project, have to invest a lot of time and effort to learn the project structure, code base, etc. before they can contribute effectively
% LLMs could potentially help with this, but it is not clear how to use them effectively in this context, especially in large software projects like CrypTool-Online
\subsection{Research Questions}
% 1. How can LLMs (AI Agents) be used to make it easier for new developers to contribute to a large software project, such as CrypTool-Online? Or: Can LLM agents make it easier for a new contributor to work in a large project?
% 2. What is possible with current online LLMs (AI Agents) and what are their limitations in this context?
% 3. What is possible with current offline LLMs (AI Agents) and what are their limitations in this context?
% 4. What is the effect for a reviewer after such a process is applied?
\subsection{Limitations for this Master's Thesis}
% Not all LLMs can be investigated, but a selection of them (e.g., Claude Code, Codex, OpenCode, Antigravity)
% Not all approaches can be investigated, but a selection of them (e.g., with/without tools, with/without skills, etc.)
% Not all apps can be implemented, but a selection of them (e.g., 2-3 apps)
% 
\subsection{Structure}

%%%
%%% Background and Related Work section
%%%
\section{Background and Related work}
% Review of the literature on LLMs, AI Agents, skills, tooling, etc.
% Review of existing approaches for using LLMs in software development, especially in large software projects
% Are there open source projects using agents?
%%%
%%% Fundamentals/Foundations/Background section
%%% LLMs, AI Agents, skills, tooling, how to make LLMs smart, etc.
%%%
\section{Foundations} % TODO: Refine subsection titles and make notes for content
\subsection{Large Language Models (LLMs)}
\subsection{Agentic AI}
\subsection{Skills and Tooling}
\subsection{Anforderungen im CTO-Entwicklungsprojekt}

% TODO: Add section to describe development of tools/skills/.md-file
\section{Development of ...} % Entwicklung eines angepassten LLMs/Agent/ vorkonfigurierten LLMs

%%%
%%% Implementation section, 
%%% how to use available Agents to solve the problem, case studies, etc.
%%%
\section{Implementation and Case Studies} %main part/hauptteil
% Verschiedene Ansätze vorbringen
% App 1: BoxCipher / CryptoBox
% App 2: Route Cipher
% App 3: Bitcoin attacks
\subsection{Experimental Setup} % (Kürzen wegen section "Development of ...")
% Describe CTO project context (repo structure, how new apps are added, etc.)
% What CTO apps are to be implemented (and why)? (How many apps should we do? 2?) State differences between apps
% How are the apps described to the AI?
% Describe the specs given to the agents in each run (one shot, only /init, tools, skills, etc. (not all combinations, but maybe 2-3 different approaches))
% Introduce selection of agents (and why)? (criteria: cost, performance, ease of use, etc.)
% Describe environment setup (i.e., terminal based, desktop app or plugin in VSC) (e.g., how to connect agents to the CTO codebase, how to evaluate their performance, etc.)
% Describe observations that will be made (e.g., (no. of) prompts, no. of iterations, time taken (elapsed time), autonomy, failures, interventions needed, etc.)
\subsection{Case Study: Claude Code} %Or: Case Study: App 1; Case Study: App 2; Case Study: App 3
% In-depth description of different approaches (e.g., with/without tools, with/without skills, etc.) 
% 1. prompt/spec given to the agent
% 2. step by step description of what the agent did
% 3. observations (where it got stuck, where intervention was needed)
% 4. results (code snippets, screenshots, etc.)
\subsection{Case Study: Codex}
% same as above (not too deep/detailed, e.g., only "dry run")
\subsection{Case Study: OpenCode}
% same as above (not too deep/detailed, e.g., only "dry run")

% Note: detailed description of the case studies, except the logs/transcripts (usb-stick, cd, appendix). Address research questions with important moments in the case studies, e.g., when the agent got stuck, when it needed intervention, etc. 

% Note: Do not investigate full matrix: N apps x M approaches x K agents, but sample from it, with a focus on the case study with Claude Code (qualitative multi-case-study). Rest could be future work. 

% Note: Use consistent structure for the case studies

%%%
%%% Evaluation section
%%% Evaluate the different approaches, compare them, discuss their strengths and weaknesses, etc.
%%%
\section{Evaluation}

%%%
%%% Discussion section
%%% Discuss the results, implications, limitations, etc.
%%%

%%% Criteria
% Define what "contribute more easily to a large software project" means (automonmy, time taken, interventions needed, etc.)
% Convention adherence/consistency (matches CTO code base / app patterns, etc.)
% Use of present code base (e.g., reuse of existing code, etc.)
% Correctness/ Maintainability / Erweiterbarkeit() of the code (does it work? does it meet the specs? etc.)
% (Maybe: Effort to learn the project structure, like /init)

% Vielleicht besser: Aspekte herausnehmen und beurteilen, welche KI besser/schlechter (oder ähnlich gut) war

%%% (Cross-Agent Comparison
% compare agents that got the same task and approach

%%% Cross-approach Comparison 
% compare different approaches (e.g., with/without tools, with/without skills, etc.) for the same agent

%%% Cross-approach Comparison (only for Claude Code)
% Within an agent, compare different approaches 
% e.g., does the setup make a difference in terms of interventions?)

%%% themes spanning across all the comparisons, general answer to the research questions

% psychological question: How does the motivation of a new developer change?

% Was ist der Unterschied mit und ohne LLM zu entwickeln?

\section{Conclusion and Future Work}
% Idea: Future work: How to prepare fixed environments for Agents (e.g., markdown files, tooling, skills) that can be used by new developers?
% Stress, that we are aware of the limits of the current LLMs

% TODO: Appendix: e.g., prompts, code snippets, etc.

%\newpage
%\listoffigures
%\listoftables
%\lstlistoflistings
%\newpage
%\printbibliography

\end{document}
